\section{A taxonomy of where calibration enters}
\label{sec:taxonomy}

The taxonomy of Figure~\ref{fig:taxonomy_main} places each work by \emph{where the calibration signal
enters the decision pipeline}. Table~\ref{tab:definitions} (Appendix~\ref{app:definitions}) gives the operational definition of each
stage as a pair: what counts, and the nearest case that does not. Five stages change how the confidence
is produced: the model's own read-off (S0), a post-hoc map (S1), prompt elicitation (S2), extra
inference (S3) and training (S4). The consumers (X) use a confidence without changing how it is
produced, and the theory and analysis works explain when training yields calibration. The last three
rows of the table define the properties recorded for every work: the output contract, the action a
confidence triggers when it is low, and whether a distribution-free guarantee is stated.


\subsection{The core: calibration in the training loss or reward (S4)}
\label{sec:core}

The \nCore core works share one property: a calibration term enters the model's own loss or reward, so
that the model itself, and not a map or a procedure around it, is changed towards calibration. We group
them by \emph{what the training signal scores}. Table~\ref{tab:compare_s4} (Appendix~\ref{app:core}) lists every core work with
its family, output contract, action, guarantee and metric.


\paragraph{T1: supervised calibration objectives (\nTOne works).} The model is fitted to confidence
targets or a calibration loss, without a reward. The oldest works regularise fine-tuning for calibration
in and out of distribution~\citep{kong2020calibrated} or teach a model to state a calibrated numeric or
verbal confidence~\citep{lin2022teaching}. Later work tunes on graded sets of correct and incorrect
answers~\citep{kapoor2024large}, distils sampling-based uncertainty into a single-pass
confidence~\citep{hager2025uncertainty,huang2025efficient}, and trains the stated confidence with a
proper scoring rule as the loss~\citep{li2025conftuner}. One line targets calibration lost elsewhere in
training: after preference alignment~\citep{xiao2025restoring} or after reinforcement
learning~\citep{xiaohu2026know}.

\paragraph{T2: reward-model calibration (\nTTwo works).} The calibrated model is a reward model, usually the one in
the training loop. Calibrated reward modelling counters the overconfidence that preference optimisation
induces~\citep{leng2024taming}; other works expose the reward model's own uncertainty as a
variance~\citep{lou2024uncertaintyaware} or calibrate process reward models by quantile regression, here to set inference compute per
instance~\citep{park2025know}. One reward-model method adds conformal
bounds~\citep{pan2026uncertaintyaware}; it is the only core work that states a distribution-free
guarantee.

\paragraph{T3: calibration rewards on the stated confidence (\nTThree works).} A reward or preference
scores the confidence the model states against the outcome. This is the largest family and, with T5, the newest:
\nTThreeRecent of its \nTThree works appeared in 2025 or 2026, \nTThreeTwentySix of them in 2026. The
rewards include the logarithmic scoring rule~\citep{baniharouni2025rewarding}, a Brier term added to the
correctness reward~\citep{damani2025beyond}, and rewards that let a simulated reader make calibrated
forecasts~\citep{band2024linguistic}. Later work asks which confidence rewards cannot be
hacked~\citep{tan2026on}, compares proper scoring rules as terminal rewards for
forecasters~\citep{turtel2026how}, separates the reasoning and calibration objectives to avoid
gradient conflict~\citep{ma2026decoupling}, and adjusts the probabilities of decision tokens
directly~\citep{yaldiz2026balancing}. Ranking-based confidence rewards~\citep{wang2026calibrationaware}
target discrimination as well as calibration.

\paragraph{T4: abstention and refusal training (\nTFour works).} The signal targets the action: the model
is trained, by supervision or by reward, to abstain when it does not know. Early work builds refusal-aware
fine-tuning data from the questions the model gets wrong~\citep{zhang2023rtuning,yang2023alignment}.
Reinforcement-learning variants use a ternary reward that scores a correct answer above an abstention
and an abstention above an error~\citep{wei2025truthrl,mohamadi2025honesty}, add clarification after a
refusal~\citep{zhai2026abstainr1}, or derive the reward from the model's knowledge
boundary~\citep{gao2026karl}. The family also carries a warning: error-penalised abstention rewards can
collapse a KL-anchored policy to refusing everything~\citep{che2026abstention}.

\paragraph{T5: process- and meaning-level confidence rewards (\nTFive works).} The confidence reward is
given per reasoning step~\citep{he2025mmboundary,wang2026process} or over answers that are equivalent in
meaning~\citep{an2025teaching,yu2026calibrating}, rather than once per final answer.

\paragraph{Theory and analysis (\nTheory works).} These works explain when training yields calibration
and are not counted in the core. Local optimality of a proper loss implies smooth
calibration~\citep{blasiok2023when}; a calibrated language model must hallucinate arbitrary facts at a rate close to the
fraction of such facts seen exactly once in training~\citep{kalai2023calibrated}, and binary-graded training and evaluation
reward guessing~\citep{kalai2025why}. Calibration that emerges from next-token training can break under
reinforcement-learning tuning~\citep{nakkiran2025trained}, and the group standard-deviation
normalisation of GRPO causes overconfidence on stochastic outcomes~\citep{bereket2025uncalibrated}.

\paragraph{Families crossed with the output contract.} The families differ in what they score; they
also differ in what form the trained confidence takes. Of the \nCore core works, \nCoreScore train a
single confidence score, \nCoreFree a free-text answer with a hedge or refusal, and \nCoreTyped the
probability of a typed decision, \coreTypedCalPhrase{} its calibration; \coreOtherPhrase{}. The typed decision appears only in T1 and T3, and abstention training (T4)
is almost entirely free text:

\begin{center}\footnotesize
\begin{tabular}{@{}lrrrrr@{}}
\toprule
Family & Score & Free text & Typed & Other & All \\
\midrule
T1 supervised calibration objectives & \nTOneScore & \nTOneFree & \nTOneTyped & \nTOneOther & \nTOne \\
T2 reward-model calibration & \nTTwoScore & \nTTwoFree & \nTTwoTyped & \nTTwoOther & \nTTwo \\
T3 calibration rewards on the stated confidence & \nTThreeScore & \nTThreeFree & \nTThreeTyped & \nTThreeOther & \nTThree \\
T4 abstention and refusal training & \nTFourScore & \nTFourFree & \nTFourTyped & \nTFourOther & \nTFour \\
T5 process- and meaning-level confidence rewards & \nTFiveScore & \nTFiveFree & \nTFiveTyped & \nTFiveOther & \nTFive \\
\midrule
All core works & \nCoreScore & \nCoreFree & \nCoreTyped & \nCoreOther & \nCore \\
\bottomrule
\end{tabular}\\[2pt]
{\scriptsize Other: a prediction set, or a contract that could not be determined from the paper.}
\end{center}

\noindent The action recorded for a
core work is most often none (\nCoreNoAction works): the trained confidence is reported with the answer,
and the rule that would act on it is left to the user. Abstention is the exception, because the T4
family trains the action itself (\nCoreAbstain core works abstain).

\subsection{The background stages and the consumers}
\label{sec:background}

% GENERATED by build/gen_handtables.py (citations resolved against the verified corpus) -- edit the generator.
\begin{table}[t]
\centering\footnotesize
\setlength{\tabcolsep}{2pt}\renewcommand{\arraystretch}{1.06}\hyphenpenalty=2000\exhyphenpenalty=2000
\caption{\textbf{What each stage yields, costs, guarantees and gets wrong.} The second column is the sense in which the
literature uses ``confidence'' at that stage. A cited phrase is a finding of the cited work; an uncited phrase follows from
the definition of the stage (Table~\ref{tab:definitions}). Only training (S4) changes the model itself; its families are
compared in Table~\ref{tab:compare_s4}.}
\label{tab:stage_limits}
\begin{tabular}{@{}>{\raggedright\arraybackslash}p{0.095\textwidth}>{\raggedright\arraybackslash}p{0.16\textwidth}>{\raggedright\arraybackslash}p{0.12\textwidth}>{\raggedright\arraybackslash}p{0.125\textwidth}>{\raggedright\arraybackslash}p{0.105\textwidth}>{\raggedright\arraybackslash}p{0.165\textwidth}>{\raggedright\arraybackslash}p{0.15\textwidth}@{}}
\toprule
\textbf{Stage} & \textbf{Confidence it yields} & \textbf{Cost: data; inference} & \textbf{Guarantee} & \textbf{Black-box use} & \textbf{Typical failure} & \textbf{Fit for a typed decision} \\
\midrule
S0 read-off & the model's own probability for an option or a true/false judgement \citep{kadavath2022language} & none; none & none & no: needs token probabilities & calibration degrades after post-training \citep{openai2023gpt4} & natural when the options are enumerated and scored \citep{kadavath2022language} \\
S1 post-hoc map & a rescaled, binned or pooled probability \citep{guo2017calibration}, a probe's probability from the hidden states \citep{azaria2023internal}, or an auxiliary model's estimate from the input and output text \citep{ulmer2024calibrating} & held-out data (labels usually; none for a fixed map); negligible & binning: distribution-free \citep{gupta2021distribution}; scaling: none & partly: a map needs the model's scores and a probe its hidden states; an auxiliary model on the text needs neither & a map fitted for one task must be refitted, or predicted by a learned model, for a new one \citep{shen2024thermometer} & good for a fixed label set \citep{guo2017calibration} \\
S2 prompt elicitation & a confidence stated in words or numbers \citep{tian2023just} & none; none if asked with the answer, one more call if asked afterwards & none & yes & stated confidences cluster near the top of the scale \citep{xiong2023can} & the confidence arrives as text and must be parsed \\
S3 sampling & agreement or semantic entropy across sampled answers \citep{wang2022selfconsistency,farquhar2024semantic} & none; several sampled answers per query & none & yes & consistent but wrong answers look confident \citep{tan2025too} & defined over free text; per-option use needs clustering \\
S3 conformal and risk control & a prediction set of options \citep{campos2024conformal}, or a threshold with bounded risk, such as an early exit \citep{schuster2022confident} & a calibration set; a threshold over one pass, or over several sampled answers for sets of generations & coverage $\ge 1-\alpha$ \citep{campos2024conformal}, or a risk bound \citep{angelopoulos2021learn} & either & the guarantee lapses when exchangeability fails \citep{ulmer2024nonexchangeable} & a set returns several options, not one decision; a risk-controlled threshold returns one \\
S4 training & a confidence the model is trained to state or to assign, per option for a typed decision \citep{yaldiz2026balancing} & outcomes or rewards; training compute & none in general; one reward-model method adds conformal bounds \citep{pan2026uncertaintyaware} & no: needs weights & a calibration and accuracy trade-off \citep{yaldiz2026balancing} & can train per-option probabilities directly \citep{yaldiz2026balancing} \\
\bottomrule
\end{tabular}
\end{table}


The stages outside the core are organised in the existing surveys~\citep{li2025survey,shorinwa2025survey}.
Table~\ref{tab:stage_limits} compares them with the core on what each yields, costs, guarantees and gets
wrong, and Table~\ref{tab:landscape} (Appendix~\ref{app:landscape}) lists all \nLandscape works.


The read-off (S0, \nSzero works) needs token probabilities; post-hoc maps (S1, \nSone works) add
temperature scaling, binning~\citep{guo2017calibration} and probes on hidden
states~\citep{azaria2023internal}; prompt elicitation (S2, \nStwo works) works from text
alone~\citep{tian2023just}. Extra inference (S3, \nSthree works) samples several answers and measures
their agreement~\citep{wang2022selfconsistency,farquhar2024semantic}, or calibrates a set or a threshold
with a conformal or risk-control guarantee~\citep{campos2024conformal}; \nGuarSthree of the \nGuar works
that state a guarantee sit here. The consumers (X, \nX works) route, defer, abstain, constrain decoding,
or study how people rely on a stated confidence.
